\subsection[The case of p=2]{The case of $\boldsymbol{p=2}$}
We need to describe the algebra $\widehat{D}(T_2(-1))$ in greater detail, paying particular attention to the element $\sigma$.

By \eqref{split} the category of DG-modules over $\widehat{D}(T_2(-1))[\theta]$ is a product of categories, $\Cc_0\times \Cc_1$, corresponding to the cases $s=0$ and $s=1$. We~will deal with both separately. More precisely, for a $\widehat{D}(T_2(-1))[\theta]$-module $V$, we decompose
\begin{gather}
V=(V_{00}\oplus V_{11})\oplus(V_{01}\oplus V_{10}).
\end{gather}

Observe that by the Corollary~\ref{gens} we have that
\begin{gather}
xx'+x'x=1+(-1)^s,
\end{gather}
so that
\begin{gather}\label{dd}
(x'x)^2=(1+(-1)^s)x'x.
\end{gather}
We see from \eqref{dd} that the minimal polynomial of $x'x$ depends only on $s$; this is exclusive to $p=2$ and makes this case tractable. Note that by \eqref{sigmaaction}:
\begin{gather}\label{sig}
\sigma|_{V_{00}}=1-x'x,\qquad \sigma|_{V_{11}}=-1+x'x,\qquad\text{and}\qquad \sigma|_{V_{01}}=\sigma|_{V_{10}}=1+x'x.
\end{gather} Let
\begin{gather}
D_s=\widehat{D}(T_2(-1))/(gg'-(-1)^s).
\end{gather}

We begin with $s=0$: The category of $D_0$-modules consists of $\mathbb{Z}/2$-graded vector spaces ($V=V_{00}\oplus V_{11}$) equipped with degree changing operators $x$ and $x'$ subject to the relation $xx'+x'x=2$. The action of $\sigma-1$ on $V_{00}$ is $-x'x$ and on $V_{11}$ it is $x'x-2$.

\begin{Lemma}
The category $\Cc_0$ consists of the mixed complexes of~{\rm \cite{kassel}}.
\end{Lemma}

 \begin{proof}
Note that $D_0/(\sigma-1)$-modules are just vector spaces. Indeed, let $y=x'/2$. For improved clarity, denote by $x_i$ and $y_i$ the action of $x$ and $y$, respectively, that originates at $V_{ii}$. After modding out by $\sigma-1$ we have by \eqref{sig} that $y_1x_0=0\implies x_1y_0=1$ and $y_0x_1=1$. So
\begin{gather}
y_0:V_{00}\simeq V_{11}: x_1.
\end{gather}
Furthermore, $x^2=y^2=0$ so that both $x_0$ and $y_1$ are the $0$ maps. Thus
\begin{gather}
V=V_{00}\oplus V_{11}\mapsto V_{00}
\end{gather}
establishes an equivalence of categories between $D_0/(\sigma-1)$-modules and $\vect$.

The action of $\sigma-1$ on $D_0$ is diagonalizable by \eqref{dd} and so by the proof of Lemma~\ref{diag} the algebras $D_0[\theta]$ ($d\theta=\sigma-1$) and $D_0/(\sigma-1)[\theta]$ ($d\theta=0$) are quasi-isomorphic. Thus, $\Cc_0$, being DG-equivalent to $D_0/(\sigma-1)[\theta]$-modules, is by the above discussion, equivalent to $k[\theta]$-modules. These are just the mixed complexes of~\cite{kassel}.
\end{proof}

 \begin{Remark}
 Note that not only does $\Cc_0$ consist of the usual mixed complexes but it also does not provide any evidence of the need for the mixed $aYD$-contramodules (see the proof above).
 \end{Remark}

Moving on to $s=1$ we find that things change for the better. Recall that $D_1$ is generated by
\begin{gather}
x,\ x',\ g
\end{gather}
subject to
\begin{gather}
x^2=x'^2=g^2-1=0, \qquad
xx'=-x'x, \qquad
gx=-xg, \qquad
gx'=-x'g.
\end{gather}
Furthermore, $\sigma-1$ acts as $x'x$ by \eqref{sig}.

\begin{Proposition}\label{notsscase}
Let $H=T_2(-1)$, then the mixed complexes in the category of stable anti-Yetter--Drinfeld contramodules are not DG-equivalent to the category of mixed anti-Yetter--Drinfeld contramodules.
\end{Proposition}
\begin{proof}
By the preceding discussion, i.e., the decomposition of the category into a product, it~suf\-fices to show that the categories $D_1[\theta]\text{-mod}$ (where $d\theta=x'x$) and $D_1/(x'x)[\theta]\text{-mod}$ (where $d\theta =0$) are not DG-equivalent.

Recall that the Hochschild cohomology $HH^i(C^\bullet)$ of a DG algebra $C^\bullet$ is an invariant of its DG category of modules~\cite{dginvar}. We~will compute $HH^{-1}$. In our simple case of a DG algebra $C=C^{-1}\stackrel{d}{\to} C^0$ concentrated in two degrees, we have
\begin{gather}
HH^{-1}(C)={\rm ker} \big(C^{-1}\stackrel{\alpha}{\to} C^0\oplus {\rm Hom}\big(C^0, C^{-1}\big)\big),
\end{gather}
where $\alpha(x)=(dx, [x,-])$.

The key observation here is that the center of $D_1$ is spanned by $1$, $xx'$, $xx'g$, while that of~$D_1/(x'x)$ is spanned by $1$. Thus, in the first case we get that $HH^{-1}$ is spanned by $xx'$ and~$xx'g$. In the second case it is spanned by $1$.
\end{proof}
